\section{讨论与限制}
\subsection{软件层边界}
框架生成的是 Android API 可观察的数据，不是 GNSS 射频信号。GNSS 轨道采用圆轨道近似，未模拟真实星历、伪距、载波相位、多路径、接收机时钟误差、导航电文和射频前端；因此不能替代硬件在环或射频仿真器。类似地，WiFi、BLE 和 CellInfo 是结构化测试对象，不会改变物理网络、无线电环境或基带状态。

\subsection{厂商 ROM 差异}
系统服务可能移动到 APEX、厂商 JAR 或独立 APK，R8 可能重命名内部方法，Binder 参数也可能随版本变化。Native Hook 还依赖函数布局、指令特征、PAC 和可执行内存权限。项目通过 Profile、多候选匹配、二进制特征检查和 fail-open 降低风险，但这不能保证未知 ROM 兼容。任何涉及 system\_server 的新适配都应先静态分析，再在独立检测器上做最小运行验证。

\subsection{数据一致性与可观测性}
环境多源一致性依赖统一时间轴。若适配器缓存滞后、监听器生命周期处理错误或某一数据面继续返回真实值，融合应用仍可能观察到残差。因此测试报告应记录配置版本、快照时间、Hook 状态、回调时间、原始/虚拟摘要和系统日志，而不能只记录 UI 文本。对于事件驱动计步器，静止状态没有真实事件，必须区分“没有事件”与“值为零”，并在需要时使用主动事件机制。

\subsection{后续方向}
后续工作包括：为 Android 14、16 及后续版本补充可复现实机证据；建立不依赖具体 ROM 偏移的 Sensor HAL 或系统扩展适配；将路线速度与统一运动状态双向约束；加入更严格的时间戳、单调性和跨源一致性指标；并把环境包回放与 CI 测试用例结合。上述方向均属于测试基础设施增强，不改变项目不针对第三方应用进程开发的架构约束。

\section{总结}
本文从问题背景、系统架构、数学模型、Android 实现和实验边界五个方面分析了 ZhangVirtualEnv。框架以 Backend Core 维护统一环境状态，以系统服务和 Framework 边界的适配器将状态转换为 Android API 数据，并通过 Profile 隔离 Android 版本与厂商 ROM 差异。路线引擎实现按时间推进的单点和折线轨迹，GNSS 引擎实现固定卫星目录、连续轨道近似、WGS-84 坐标变换、ENU 几何和连续 C/N0；运动引擎以统一步态相位生成步数、加速度、重力和陀螺仪数据；CellInfo、WiFi、BLE 和 SIM 适配器则通过版本兼容构造器和 Binder 服务边界提供结构化测试输入。

源码和逆向证据同时表明，数据路径的“可调用”并不等于“已送达”。尤其是 Android 15 Oplus SensorService 的 Runtime Sensor 通道存在真实传感器句柄边界，Native 分发汇聚点虽是合理的系统级研究方向，也必须由 VirEnvDetector、Hook 计数、logcat 与系统调用共同证明。因此，本文最终将 ZhangVirtualEnv 定位为可重复的 Android 系统环境参数测试与兼容性验证框架，而非真实物理环境或射频信号仿真器。
